% META: {"id": "1SPE-TRIGO-EX-050", "chapitre": "1SPE-TRIGONOMETRIE", "type_objet": "exercice", "capacites": ["1SPE-TRIGONOMETRIE-C1","1SPE-TRIGONOMETRIE-C2","1SPE-TRIGONOMETRIE-C3","1SPE-TRIGONOMETRIE-C4","1SPE-TRIGONOMETRIE-C5"], "capacites_codes": ["C1","C2","C3","C4","C5"], "methodes": ["M1","M2","M3","M4","M5"], "parcours": 3, "competences": ["chercher", "modeliser", "raisonner", "communiquer"], "duree_min": 20, "mode_creation": "ex_nihilo", "sources_inspiration": [], "parametres_sympy": {}, "fichier_tex": "chapitres/1SPE-TRIGONOMETRIE/exercices/1SPE-TRIGO-EX-050.tex", "corrige_tex": "chapitres/1SPE-TRIGONOMETRIE/corriges/1SPE-TRIGO-CO-050.tex", "status": "generated"}
% BEGIN-VERIFY
% from sympy import *
% t = symbols('t', real=True)
% # h(t) = 10 + 8*sin(pi*t/6 - pi/2) = 10 - 8*cos(pi*t/6)
% h = 10 - 8*cos(pi*t/6)
% assert h.subs(t, 0) == 2
% assert h.subs(t, 6) == 18
% assert h.subs(t, 12) == 2
% # h(t) = 14 <=> cos(pi*t/6) = -1/2 <=> pi*t/6 = 2*pi/3 + 2k*pi or -2*pi/3 + 2k*pi
% # t = 4 + 12k or t = -4 + 12k => on [0;12]: t=4 and t=8
% assert simplify(h.subs(t, 4) - 14) == 0
% assert simplify(h.subs(t, 8) - 14) == 0
% END-VERIFY
\begin{exercice}{1SPE-TRIGO-EX-050}{3}{20}
\textbf{Probleme de synthese — Marees.}

La hauteur d'eau (en metres) dans un port est modelisee par
\[
  h(t) = 10 - 8\cos\!\left(\frac{\pi t}{6}\right),
\]
ou $t$ est le temps en heures depuis la maree basse.
\begin{enumerate}
  \item Quelle est la hauteur d'eau a maree basse ($t = 0$) ? A maree haute ?
  \item Quelle est la periode du phenomene (en heures) ?
  \item Dresser le tableau de variations de $h$ sur une periode.
  \item Un bateau a besoin d'au moins $14$~m d'eau pour entrer dans le port. Resoudre $h(t) \geqslant 14$ sur $[0\,;\,12]$ et en deduire la plage horaire pendant laquelle le bateau peut entrer.
  \item Convertir les bornes de cette plage en heures et minutes.
\end{enumerate}
\end{exercice}
